\paragraph{\trDoubleBranchedCover{}。}\label{sub:double-branched-cover}
取定一个\trExact\trLefschetzFibration{}$\pi\colon E^{2n+2}\to S$，其中 $S=D$ 是 $\CC$ 中的单位\trDisk（带标准\trComplexStructure）。取一个小 $\delta>0$，把半径 $1-4\delta$ 的\trDisk{}记作 $S^{\mathrm{shrink}}\subset S$，并把它在\trTotalSpace{}中的\trPreimage{}记作 $E^{\mathrm{shrink}}\subset E$。我们假定 $\pi$ 的全部\trCriticalValue{}都落在 $S^{\mathrm{shrink}}$ 的内部，并且进一步假定：$E$ 上的负\trLiouville{}流在某个正的时刻会把整个 $E$ 推入 $E^{\mathrm{shrink}}$ 之内。对充分小的 $\delta$，这两个条件当然都得到满足。我们将取 $*=1-3\delta\in\mathrm{int}(S)$ 作为基点，并取该点上的\trFiber{}$M=E_*$。每当需要选取\trBasePath{}$\chi$ 时，我们都取从 $*$ 到 $1\in\partial S$ 的直线。
设 $\widetilde{S}$ 是 $S$ 在 $*$ 处分歧的\trDoubleCover，类似地，$\widetilde{E}$ 是 $E$ 沿 $M=E_*$ 分歧的\trDoubleCover。当然，$\widetilde{S}$ 仍然\trIsomorphic{}于 $D$，但我们宁愿不作这一认同，而是写成\footnote{译注：原书此处印作 $ty^2+*=\pi(x)$；疑应为 $y^2+*=\pi(x)$（依上行 $\widetilde{S}$ 的定义及\cref{sub:family-double-covers}中家族版定义方程 $y^2+w=\pi(x)$，字母 $t$ 在此未定义，4x 放大实证确有此 $t$）；此处径改。——译者注}
\begin{equation}\label{eq:double-cover-models}
\begin{gathered}
\widetilde{S}=\{y\in\CC:\ y^2+*\in S\},\\
\widetilde{E}=\{(y,x)\in\widetilde{S}\times E:\ y^2+*=\pi(x)\}.
\end{gathered}
\end{equation}
$\widetilde{E}$ 自带一个诱导的\trAlmostComplex{}结构（它是乘积 $\widetilde{S}\times E$ 的一个\trAlmostComplex\trSubmanifold），这使得 $\partial\widetilde{E}$ 是弱\trConvex{}的。相对于这一结构，我们有下列\trPseudoholomorphic{}映射
\begin{equation*}
\begin{array}{ll}
\left.\begin{aligned}
\iota_{\widetilde{S}}(y)&=-y\\
\iota_{\widetilde{E}}(y,x)&=(-y,x)
\end{aligned}\right\} & \text{\trCovering\trInvolution；}\\[14pt]
\left.\begin{aligned}
\beta_{\widetilde{S}}(y)&=y^2+*\colon\widetilde{S}\longrightarrow S\\
\beta_{\widetilde{E}}(y,x)&=x\colon\widetilde{E}\longrightarrow E
\end{aligned}\right\} & \text{\trQuotient{}映射；}\\[14pt]
\tilde{\pi}(y,x)=y\colon\widetilde{E}\longrightarrow\widetilde{S}, & \text{$\pi$ 的\trLift。}
\end{array}
\end{equation*}
$\tilde{\pi}$ 的\trCriticalPoint{}是 $\pi$ 的\trCriticalPoint{}的\trPreimage；由于局部结构保持不变，它们依然\trNondegenerate。
为完成这一构造，我们还需引进适当的\trExactSymplectic{}结构。回忆作为我们\trFibration{}定义数据的一部分，底空间 $S$ 带有某个 $\theta_S\in\Omega^1(S)$ 与 $\phi_S=d\theta_S\in\Omega^2(S)$。由 \eqref{eq:double-cover-models} 可以看出，$\widetilde{S}$ 包含原点周围半径 $\sqrt{2\delta}$ 的\trDisk。取 $\kappa$ 为\trSupport{}落在半径 $\sqrt{\delta}$ 的更小\trDisk{}之内的一个 1-形式，使之成为 $\iota_{\widetilde{S}}$-不变的，并且使 $d\kappa$ 在原点处为正。我们愿意把 $\kappa$ 想成小的，因此保留必要时把最初选取乘以某个正的小常数的权利。令
\begin{equation}\label{eq:shrunk-base-forms}
\theta_{\widetilde{S}}=\beta_{\widetilde{S}}^*\theta_S+\kappa,\qquad
\phi_{\widetilde{S}}=d\theta_{\widetilde{S}}=\beta_{\widetilde{S}}^*\phi_S+d\kappa.
\end{equation}
由于 $d\kappa$ 小而在原点处为正，$\phi_{\widetilde{S}}$ 处处为正；又由于 $\kappa$ 的\trSupport{}在内部，$\theta_{\widetilde{S}}$ 沿边界为正。这些正是让 \eqref{eq:shrunk-base-forms} 把 $\widetilde{S}$ 做成一个带边\trExactSymplectic\trManifold{}的条件。与此平行地，我们为 $\widetilde{E}$ 配备
\begin{equation}\label{eq:total-forms-tilde-e}
\theta_{\widetilde{E}}=\beta_{\widetilde{E}}^*\theta_E+\tilde{\pi}^*\kappa,\qquad
\phi_{\widetilde{E}}=d\theta_{\widetilde{E}}=\beta_{\widetilde{E}}^*\phi_E+\tilde{\pi}^*d\kappa.
\end{equation}
第一项 $\beta_{\widetilde{E}}^*\phi_E$ 只在分歧\trFiber{}$\{y=0\}$ 的法方向上不成辛，于是按与前面相同的论证，加上一个充分小的 $\tilde{\pi}^*d\kappa$ 项便使它处处辛。注意，$\tilde{\pi}$ 的\trRegular{}部分上与之相伴的\trSymplectic\trConnection{}是 $\pi$ 的那个的\trPullback，因为 \eqref{eq:total-forms-tilde-e} 中的第二项来自底空间。这表明 $\partial^h\widetilde{E}$ 是\trHorizontal{}的，而且（用与\cref{thm:fibre-exact-symplectic}相同的论证）负\trLiouville{}流沿边界的这一部分指向内侧。在 $\partial^v\widetilde{E}$ 附近，负\trLiouville\trVectorField{}是 $E$ 上那个的\trPullback，因而同样指向内侧。有了这些，我们便已验证 $\tilde{\pi}\colon\widetilde{E}\to\widetilde{S}$ 仍是一个\trExact\trLefschetzFibration。
在接下来的讨论中，$\widetilde{E}$ 中有两类\trExact\trLagrange\trSubmanifold{}将占据突出的位置。我们把它们记作 (U)（非分歧，unbranched）与 (B)（分歧，branched）：
\begin{itemize}
\item \emph{(U) 型。}$\beta_{\widetilde{S}}$ 限制到 $S^{\mathrm{shrink}}$ 上是一个非\trBranchedCover{}的\trCovering，因而是平凡的，这意味着\trTotalSpace{}由底空间的两份拷贝组成。$\beta_{\widetilde{E}}$ 限制到 $E^{\mathrm{shrink}}$ 上的情形亦同。我们写成
\begin{equation}\label{eq:shrink-double-components}
\begin{gathered}
\widetilde{S}^{\mathrm{shrink}}=\beta_{\widetilde{S}}^{-1}(S^{\mathrm{shrink}})=\widetilde{S}^{\mathrm{shrink}}_+\sqcup\widetilde{S}^{\mathrm{shrink}}_-,\\
\widetilde{E}^{\mathrm{shrink}}=\beta_{\widetilde{E}}^{-1}(E^{\mathrm{shrink}})=\widetilde{E}^{\mathrm{shrink}}_+\sqcup\widetilde{E}^{\mathrm{shrink}}_-.
\end{gathered}
\end{equation}
更具体地说，满足 $|y^2+*|\le 1-4\delta$ 的点 $y$ 必然有 $\mathrm{im}(y)\neq 0$，而该虚部的符号区分开 \eqref{eq:shrink-double-components} 的两个\trConnectedComponent。这样的 $y$ 还满足 $|y|\ge\sqrt{\delta}$，这意味着：在每个\textbf{\trDoubleBranchedCover}（double branched cover）的限制版本上，\eqref{eq:shrunk-base-forms}与\eqref{eq:total-forms-tilde-e} 中的 $\kappa$ 项都成为零。一个显然的推论是：若 $L$ 是 $E$ 的一个含于 $E^{\mathrm{shrink}}$ 之内的\trExact\trLagrange\trSubmanifold，则它在 $\widetilde{E}$ 中的\trPreimage{}由 $L$ 的两份拷贝组成，而它们仍是\trExact{}的\trLagrange\trSubmanifold。仿照上面的做法，我们使用记号
\begin{equation*}
\widetilde{L}=\beta_{\widetilde{E}}^{-1}(L)=\widetilde{L}_+\sqcup\widetilde{L}_-.
\end{equation*}
\item \emph{(B) 型。}设 $\gamma\colon[0;1]\to S$ 是一条满足 $\gamma(0)=*$ 的\trVanishingPath，并设 $\Delta_\gamma\subset E$ 是它的\trLefschetzThimble。那么
\begin{equation}\label{eq:branched-preimage-thimble}
\widetilde{\Delta}_\gamma=\beta_{\widetilde{E}}^{-1}(\Delta_\gamma)
\end{equation}
是 $\widetilde{E}$ 中的一个\trSmooth\trLagrange\trSphere。看出这一点的一种方式，是把它表示成一个\trMatchingCycle，如下。在相差\trOrientation{}反转的意义下，存在唯一的\trSmooth{}路径 $\tilde{\gamma}\colon[-1;1]\to\widetilde{S}$ 使得 $\beta_{\widetilde{S}}(\tilde{\gamma}(t))=\gamma(t^2)$。$\tilde{\gamma}$ 的两个端点是\trCriticalValue{}$\gamma(1)$ 的\trPreimage，因此其两半 $\tilde{\gamma}_\pm(t)=\tilde{\gamma}(\pm t)\colon[0;1]\to\widetilde{S}$ 都是 $\tilde{\pi}$ 的\trVanishingPath。如上文所指出的，\trSymplectic\trConnection{}是 $\pi$ 的那个的\trPullback，由此可知 $\tilde{\gamma}_\pm$ 的每个\trLefschetzThimble{}都投射到 $\gamma$ 的\trLefschetzThimble{}之上。这意味着 $\tilde{\gamma}_\pm$ 的诸\trVanishingCycle{}相重合——都等于 $M$ 内 $\gamma$ 的\trVanishingCycle——而且 \eqref{eq:branched-preimage-thimble} 恰是 $\tilde{\gamma}$ 按朴素意义所取的\trMatchingCycle；参见\cref{sub:matching-cycles}的开头。
\end{itemize}
不用 $*=1-3\delta$，人们也可以取另一个\trRegularValue{}$w\in S\setminus\mathrm{Critv}(\pi)$ 作为\trBranchPoint。证明\trDoubleBranchedCover{}是\trLefschetzFibration{}的技术性细节要求 $|w|\le 1-2\delta$。对 (U) 型\trLagrange\trSubmanifold{}的构造，则不得不假定 $L$ 落在区域 $|\pi(x)|\le|w|-\delta$ 之内。此外，$\widetilde{L}_+$ 与 $\widetilde{L}_-$ 之间的区分还依赖于 $\sqrt{w}$ 的选取。这当然意味着：如果让 $w$ 绕、比如说、半径 $1-3\delta$ 的\trCircle{}移动，两个\trConnectedComponent{}便会互换。我们将在下文把这一点更严格化，见\cref{thm:monodromy-involution-agreement}。
